% ============================================================================
%  Subdocumento 12. Equipo de trabajo, subcontrataciones y alianzas (D4-153).
%  El equipo nominado es el del Subdocumento 1. La dotación sale del
%  Formulario T-15 y los porcentajes de los proveedores del Sobre N.º 3.
% ============================================================================

\begin{resumenApertura}
Este subdocumento presenta quién hace el proyecto y quién responde por cada parte. El equipo clave son
los seis roles que exige el \art{34.1}{22}, nominados y con su dedicación en el Formulario T-8. La
dotación sube de 14 personas en los meses 1 y 2 a 37 en los meses 13 a 17 y baja a 16 en la operación,
igual que en el Formulario T-15. Lo que no hace audIT lo hacen proveedores y socios que suman cerca del
15,6 \% del valor del contrato, bajo el tope del 40 \% del \art{73.1}{38}.
\begin{recibe}
  \item Una estructura del proyecto con un responsable por cada entregable y una matriz de asignación.
  \item Un equipo clave nominado, con certificaciones, dedicación y meses de participación en el Formulario T-8.
  \item Una curva de dotación por fase que coincide con la nivelación de recursos del Formulario T-15.
  \item La lista de subcontratistas y socios con su rol, su peso en el valor y su régimen de control.
\end{recibe}
\end{resumenApertura}

El equipo es el que declara el Subdocumento 1: 22 profesionales de planta de audIT, de los cuales 21
trabajan en el proyecto, y 16 personas de refuerzo que audIT incorpora para el despliegue y la segunda
etapa \form*{T-15}{65}. La sección 12.1 ordena el proyecto, la 12.2 presenta el equipo clave, la 12.3 la
dotación por fase, la 12.4 lo que audIT hace y lo que compra, la 12.5 los subcontratistas y socios, y la
12.6 cómo se conserva el conocimiento cuando alguien se va.

% ============================================================================
\section{Estructura organizacional del proyecto}\label{sec:12-estructura-organizacional-del-pr}
% ============================================================================

El proyecto tiene cuatro líneas bajo el Jefe de Proyecto, como muestra la \figref{12-organigrama}. La
dirección técnica diseña y construye. El aseguramiento vigila la seguridad y la calidad con
independencia de quien construye. La implantación lleva la solución a los terminales, a los camiones y
a las personas. La operación sostiene el servicio desde que entra en producción. Las decisiones que
cruzan las líneas suben a los comités del \art{71}{37}: el Comité de Proyecto cada quince días, el de
Arquitectura cada mes, el Ejecutivo cada mes y, desde el mes 13, el de Operación.

\begin{figuraNativa}[!ht]{Estructura del proyecto: el Jefe de Proyecto y sus cuatro líneas, con los roles de cada una}{12-organigrama}
\begin{tikzpicture}[x=1mm,y=1mm,
  caja/.style={draw=audit-marino,line width=0.6pt,fill=white,rounded corners=2pt,align=center,
               text width=34.5mm,inner sep=2pt,anchor=north,
               font=\sffamily\fontsize{9.5bp}{11bp}\selectfont\color{audit-tinta}},
  jefe/.style={caja,fill=audit-gris-claro,text width=60mm},
  linea/.style={line width=0.6pt,draw=audit-marino}]
  \node[jefe] (jp) at (55,0) {\textbf{Jefe de Proyecto}\\ Comités del Artículo 71 y oficina de proyecto};
  \node[caja] (t) at (0,-24) {\textbf{Dirección técnica}\\ Arquitecto de Solución\\ Líder de Datos\\ Líder de Desarrollo\\ Líder de Integración};
  \node[caja] (a) at (37,-24) {\textbf{Aseguramiento}\\ Encargado de Seguridad\\ de la Información\\ Líder de Calidad};
  \node[caja] (i) at (74,-24) {\textbf{Implantación}\\ Líder de Implantación\\ y Gestión del Cambio\\ Líder Funcional\\ Analistas en terreno};
  \node[caja] (o) at (111,-24) {\textbf{Operación}\\ Líder de Operación\\ y gerente de servicio\\ Mesa, plataforma\\ y equipo a bordo};
  \draw[linea] (jp.south) -- (55,-17);
  \draw[linea] (0,-17) -- (111,-17);
  \foreach \x in {0,37,74,111} { \draw[linea] (\x,-17) -- (\x,-24); }
\end{tikzpicture}
\end{figuraNativa}

El Jefe de Proyecto responde ante el mandante por el contrato completo. El Arquitecto de Solución decide
la arquitectura y lleva su registro de decisiones. El Encargado de Seguridad de la Información y el Líder
de Calidad dependen del Jefe de Proyecto y no de la dirección técnica, para que puedan detener una entrega
que no cumple. El Líder de Operación es además el gerente de servicio dedicado que pide
\transv{RT-21.02}{35}. La \tabref{12-raci} agrupa los paquetes del Formulario T-14 y marca quién
responde por cada grupo (A), quién lo ejecuta (R), a quién se consulta (C) y a quién se informa (I).

\begin{tabla}[ancho=completo, fuente=condensada]{Matriz de asignación por grupo de paquetes de la EDT}{12-raci}{X *{10}{C{8.2mm}}}
Paquetes de la EDT & JP & AS & ESI & LD & LDS & LC & LO & LIN & LIM & LF \\ \midrule
1.1 y 1.2 Dirección y comunicaciones & A & C & C & I & I & C & C & I & C & C \\
2.1 Alcance y trazabilidad & C & C & I & C & C & C & I & C & I & A \\
2.4 Arquitectura, seguridad y modelo de datos & I & A & R & R & C & C & C & C & I & I \\
3.1, 3.3 y 3.4 Ambientes, trazas y sala & I & C & C & I & C & C & A & I & I & I \\
3.2 Identidad y cifrado & I & C & A & C & C & C & R & I & I & I \\
4.2 y 11.1 Montaje a bordo y capacitación & I & I & I & I & I & C & C & C & A & C \\
4.3 y 6.1 a 6.5 Integraciones y retiro del sistema de 2013 & I & C & C & C & R & C & C & A & I & I \\
5.1 a 5.7 y 8.1 a 8.4 Contextos y portales & I & C & C & C & A & C & I & C & I & C \\
7.1 a 7.3 y 8.5 Migración y analítica & I & C & C & A & R & C & I & C & I & C \\
9.1 a 9.3 Adhesión de transportistas & C & I & C & I & I & I & I & I & C & A \\
10.1 a 10.3 Pruebas & I & C & C & C & R & A & C & R & C & C \\
11.2 y 11.3 Marchas blancas & C & C & C & C & R & C & R & R & A & C \\
11.4 y 13.1 a 13.3 Estabilización y operación & I & C & C & C & C & C & A & C & I & I \\
13.4 Mantención evolutiva & I & C & C & C & A & C & C & C & I & I \\
12.1 a 12.3 y 12.5 Innovaciones & C & A & C & C & R & C & C & C & I & C \\
\end{tabla}
\notaTabla{JP, Jefe de Proyecto. AS, Arquitecto de Solución. ESI, Encargado de Seguridad de la Información. LD, Líder de Datos. LDS, Líder de Desarrollo. LC, Líder de Calidad y Pruebas. LO, Líder de Operación. LIN, Líder de Integración. LIM, Líder de Implantación y Gestión del Cambio. LF, Líder Funcional. Fuente: elaboración propia sobre el Formulario T-14.}

Cada grupo tiene una sola A, y es el mismo rol responsable que el Formulario T-14 asigna a sus paquetes.
Los paquetes 1.3, 2.2, 2.3, 2.5, 4.1 y 12.4 no aparecen en la tabla porque su responsable es un rol que
ya figura en ella o un ingeniero de especialidad, y están en el Formulario T-14. La seguridad y la calidad
tienen su propio responsable y el desarrollo solo ejecuta en esos grupos, que es la separación que exige
el aseguramiento independiente del Subdocumento 9.

% ============================================================================
\section{Equipo clave}\label{sec:12-equipo-clave}
% ============================================================================

El equipo clave son los seis roles habilitantes del \art{34.1}{22}. El Subdocumento 1 los nomina y deja a
este subdocumento su trayectoria, sus certificaciones y su dedicación. Los seis son profesionales de
planta de audIT con entre 9 y 14 años de experiencia en el trabajo que les toca en este contrato.

El Jefe de Proyecto es Alejandro Hermosilla Díaz, certificado PMP, con 14 años dirigiendo proyectos de
modernización de plataformas digitales y sistemas de misión crítica en logística y transporte. Conduce
la planificación, los comités del \art{71}{37} y la relación con la contraparte técnica. El Arquitecto
de Solución es Esteban Valenzuela Lagos, con las certificaciones Microsoft Azure Solutions Architect
Expert y TOGAF 9.2 y 12 años en arquitecturas de nube híbrida, microservicios y sistemas de alta
concurrencia. Decide la arquitectura de los Subdocumentos 4 y 5 y lleva el Comité de Arquitectura junto
al Encargado de Seguridad.

El Encargado de Seguridad de la Información es Mauricio Arancibia Toledo, con CISSP, CISM y auditor
líder ISO/IEC 27001, y 11 años en ciberseguridad industrial, criptografía y protección de datos
personales conforme a la Ley N.º~21.719. Responde por el plan de seguridad, el cifrado de campo y la
certificación ISO/IEC 27001 de audIT. El Líder de Datos es Rodrigo Sanhueza Parra, con CDMP y Azure
Data Engineer Associate, y 9 años en datos distribuidos, flujos de eventos en tiempo real y bases
relacionales y de series de tiempo. Responde por el modelo de datos, la migración de las cerca de 6.000
vigencias y el repositorio analítico.

La Líder de Calidad es Claudia Navarrete Rivas, con ISTQB Full Advanced Level como Test Manager y
Technical Test Analyst, y 10 años dirigiendo compuertas de calidad y pruebas automatizadas de regresión,
carga y estrés. Responde por el plan de calidad del Subdocumento 9 y por la aceptación de cada hito. El
Líder de Operación es Felipe Morales Cárdenas, con CKA y Red Hat Certified Engineer, y 10 años en
orquestación de Kubernetes, telemetría industrial, redes vehiculares SAE J1939 y soporte 24/7 en
terreno. Dirige la plataforma, la mesa de servicio y el ciclo de vida del equipo a bordo, y en la
operación es el gerente de servicio.

Los seis trabajan a tiempo completo en los meses 1 a 20. En la operación, según el Subdocumento 1, el
Líder de Operación sigue a tiempo completo, el Jefe de Proyecto, el Encargado de Seguridad y el Líder de
Datos quedan con la mitad de su tiempo, y el Arquitecto de Solución y la Líder de Calidad con un cuarto.
Eso suma tres jornadas completas, 480 de las 2.560 horas mensuales de la operación, y esas horas están
dentro de las 16 personas del Formulario T-15, que es también como las cuenta el Sobre N.º~3. El Líder
de Datos y la Líder de Calidad las dedican sobre todo al paquete 13.4, donde están los cierres de los
meses 21 a 24 (Subdocumento 11). El Formulario T-8 trae cada rol con su nombre, sus certificaciones, su
dedicación y sus meses, y los currículos y las cartas de compromiso van en el Sobre N.º~1. Los demás
roles de la estructura salen de la planta o del refuerzo que declara el Formulario T-15, y se nominan en
el plan de dirección del mes 1.

% ============================================================================
\section{Curva de dotación por fase}\label{sec:12-curva-de-dotacion-por-fase}
% ============================================================================

La dotación sigue la nivelación de recursos del Formulario T-15 mes a mes, como muestra la
\figref{12-dotacion}. El proyecto parte con 14 personas en los meses 1 y 2, que hacen el levantamiento y
el plan, y sube a 24 en los meses 3 a 5 con la arquitectura y los ambientes. Llega a 30 en los meses 6 a
9, cuando empiezan la construcción y el montaje a bordo, y a 33 en los meses 10 a 12, con la
capacitación. El peak es de 37 personas en los meses 13 a 17, cuando la marcha blanca y la
estabilización de la Etapa 1 conviven con la construcción de la Etapa 2. De esas 37, 21 son de planta y
16 de refuerzo. Después baja a 30 en el mes 18 y a 26 en los meses 19 y 20, y la operación sigue con 16
personas desde el mes 21 hasta el 56. La implementación suma 93.760 horas y la operación 92.160.

\begin{figuraNativa}[!ht]{Personas por mes según el Formulario T-15, con la parte que corresponde al equipo clave}{12-dotacion}
\begin{tikzpicture}[x=2.3mm,y=1.15mm,
  t/.style={font=\sffamily\fontsize{9.5bp}{11bp}\selectfont\color{audit-tinta}},
  clave/.style={fill=audit-marino!80,draw=white,line width=0.4pt},
  resto/.style={fill=audit-turquesa!35,draw=white,line width=0.4pt}]
  \foreach \a/\b/\n in {0/2/14,2/5/24,5/9/30,9/12/33,12/17/37,17/18/30,18/20/26} {
    \fill[clave] (\a,0) rectangle (\b,6);
    \fill[resto] (\a,6) rectangle (\b,\n);
    \node[t,anchor=south] at ({(\a+\b)/2},\n) {\n};
  }
  \fill[clave] (20,0) rectangle (56,3);
  \fill[resto] (20,3) rectangle (56,16);
  \node[t,anchor=south] at (38,16) {16};
  \draw[audit-filete,line width=0.6pt] (0,0) -- (56,0);
  \foreach \m/\x in {1/0.5,13/12.5,21/20.5,56/55.5} {
    \draw[audit-filete] (\x,0) -- (\x,-1.2);
    \node[t,anchor=north] at (\x,-1.6) {\m};
  }
  \node[t,anchor=north] at (28,-6.2) {Mes del contrato};
  \fill[clave] (30,33) rectangle (31.5,35.5);
  \node[t,anchor=west] at (31.8,34.25) {Equipo clave, 6 y luego 3 jornadas};
  \fill[resto] (30,28.5) rectangle (31.5,31);
  \node[t,anchor=west] at (31.8,29.75) {Resto del equipo};
\end{tikzpicture}
\end{figuraNativa}

La forma de la curva tiene una razón. El peak contiene los meses 13 a 15, en que el Subdocumento 7
divide el equipo en dos escuadras, una que estabiliza lo que ya opera y otra que construye lo que viene,
y por eso no se puede aplanar sin atrasar la Etapa 2. Las 16 personas de refuerzo son el Líder
Funcional, el Líder de Integración, el Líder de Implantación y Gestión del Cambio, siete analistas de
implantación en terreno, dos ingenieros SRE, dos desarrolladores, un ingeniero de datos y un analista
de pruebas. Las 16 personas de la operación se reparten en la mesa de servicio, la plataforma, el
equipo a bordo y la mantención evolutiva, con el dimensionamiento de los Subdocumentos 10 y 11.

La curva también se comprobó contra el trabajo por tarea, como pidió la revisión del Informe 2. El mes 9 es el más cargado de construcción. La estimación por actividad del
software, los datos y las integraciones pide ese mes 3.354 horas de arquitectura, desarrollo, interfaces,
seguridad, datos, pruebas y plataforma, unas 21 jornadas completas. Sumados el montaje a bordo, las
pruebas transversales, la dirección y la preparación de las innovaciones, la demanda del mes llega a
cerca de 4.283 horas, y de las 4.800 que dan las 30 personas del Formulario T-15 quedan unas 517 para
las cargas menores que faltan por detallar. Por perfil, seguridad y plataforma quedan justo en su
capacidad ese mes, con unas 10 y 21 horas por sobre una jornada, y esa diferencia la toman otros
ingenieros de planta con el perfil compatible. Las cuentas
de usuario del repositorio siguen la misma curva: 33 en la primera etapa, 37 en el peak y 16 desde el
mes 25.

% ============================================================================
\section{Decisiones de hacer o comprar}\label{sec:12-decisiones-de-hacer-o-comprar}
% ============================================================================

La regla es construir lo que es propio del negocio de Curimón, comprar lo que el mercado ya hace mejor y
contratar lo que exige independencia o una especialidad que no conviene mantener en planta. Cada compra
se compara por su costo total de 56 meses, que incluye alta, canon, consumo, alojamiento, actualización,
soporte y salida, y por compatibilidad, permisos sobre los datos, ubicación, portabilidad y derecho de
uso. Un precio bajo no basta para decidir.

audIT construye el software de los seis contextos, las reglas de asignación y liquidación y las
interfaces web y móviles. La alternativa era comprar un sistema de gestión de transporte y
parametrizarlo, pero ninguno acredita cubrir a la vez la verificación de jornada, el documento de
transporte y la convivencia con el sistema de 2013 que piden los Subdocumentos 3 y 4. También construye
los adaptadores hacia el sistema contable, la telemetría de fábrica y las plataformas de posicionamiento,
sobre contratos autorizados y sin tocar sus bases, porque el sistema contable conserva la emisión
tributaria y cada dato tiene su dueño. La migración se construye como un proceso reproducible con
revisión funcional, porque las cerca de 6.000 vigencias se verifican una por una contra su respaldo. Lo
que se construye es del mandante desde su creación \art*{84.1}{43}. audIT tiene para eso siete
ingenieros de software y cinco de firmware y hardware en su planta, certificación ISO 9001 vigente y los
tres proyectos de misión crítica del Formulario T-6 (Subdocumento 1).

audIT compra como servicio lo que no es propio del negocio. Las bases de datos, la caché, los eventos, la
identidad y la custodia de claves son servicios administrados de Azure, porque operar clústeres propios
recargaría a un área de tecnología de nueve personas \caso*{numeral 7.4}{16} con parches, alta
disponibilidad y respaldos. audIT es socio Microsoft CSP, que es la condición de socio del proveedor de
nube que pide el \art{34.1}{22}. Las bibliotecas y herramientas abiertas se usan con su inventario de
componentes y no se reimplementan. GitLab se contrata en su edición Premium autoadministrada, dentro del
perímetro del proyecto, en vez de la gratuita o la de software como servicio. SonarQube se contrata en su
edición Developer porque la gratuita no analiza Dart ni las ramas de integración que usa el
Subdocumento 6.

El hardware se compra a fabricantes y no se fabrica. Los equipos de campo y de red los compra el
mandante con la especificación del Formulario T-11, y los servidores del nodo de continuidad de San
Bernardo los provee audIT.
Los servicios que la solución consume, como Azure, el circuito ExpressRoute, los planes móviles y el
tiempo de aire satelital, los contrata y paga audIT, con la suscripción de Azure a nombre del mandante. La
instalación a bordo no se subcontrata: la hace el personal de audIT que integró el equipo en su
laboratorio.

audIT contrata lo que pide independencia o una especialidad puntual. Las pruebas de penetración y la
custodia de las fuentes las hace un tercero porque su valor está en no depender de quien construye. Un
laboratorio acreditado ensaya el equipo a bordo contra polvo, agua y vibración entre los meses 4 y 7, y
el banco de pruebas del bus CAN se arrienda entre los meses 4 y 9. La plataforma de aprendizaje en línea
y sus contenidos se contratan para la capacitación de la Etapa 1, y los talleres presenciales en los
cinco terminales, la capacitación de la Etapa 2 y la reinducción anual de la operación los dicta un
organismo de capacitación. Los talleres de ergonomía de la innovación 5 también son externos.

% ============================================================================
\section{Subcontratistas y socios}\label{sec:12-subcontratistas-y-socios}
% ============================================================================

El Formulario T-8 lista cada proveedor y socio con el servicio que da y su peso en el valor total que
factura el contrato. El mayor es Microsoft, con cerca del 8,8 \%, por el consumo de Azure y Entra ID.
Le siguen Ground Control, con cerca del 3,2 \% por el tiempo de aire satelital, y los fabricantes y
mantenedores del hardware, con cerca del 1,5 \%. El resto pesa menos del 1 \% cada uno. En total suman
cerca del 15,6 \%, bajo el tope del 40 \% del \art{73.1}{38}. Los porcentajes salen de las partidas del
Sobre N.º~3, que es donde están los montos. Los cuatro centros técnicos aliados de la red regional del
Subdocumento 1 también se declaran en el formulario, aunque no tienen una partida propia en el Sobre
N.º~3.

Ningún proveedor opera el núcleo del negocio, como prohíbe el \art{73.2}{38}: las aplicaciones las
construye y opera audIT. Los datos sí quedan en Azure, y por eso la suscripción está a nombre del
mandante, que autoriza expresamente al proveedor de nube. Ese es también el proveedor que más pesa, y el
riesgo de quedar atado a él lo acota la estrategia de reversibilidad del Subdocumento 4, con la
infraestructura como código y los datos exportables en formatos abiertos.

El régimen de control tiene cinco partes. Primero, cada subcontratación tiene la autorización previa y
escrita del mandante \art*{73.1}{38}, y también la tiene el reemplazo de un proveedor, que debe
mantener o mejorar las condiciones del anterior. Segundo, todo proveedor con acceso a datos del mandante
queda declarado, firma un acuerdo de tratamiento de datos conforme a la Ley N.º~21.719 y queda sujeto a
auditoría \art*{73.4}{38}. Tercero, cada contrato traslada al proveedor los plazos que audIT necesita para
cumplir los niveles de servicio, restados desde el \art{78.2}{40} (Subdocumento 10). Cuarto, el
proveedor cumple los mismos estándares de seguridad, confidencialidad, calidad y cumplimiento laboral
que audIT, que le pide sus certificados laborales y previsionales al contratarlo y en cada renovación
\art*{73.3}{38}. Quinto, el Comité de Proyecto revisa a los proveedores cada quince días durante la
implementación y el Comité de Operación cada mes desde el mes 13, con el informe de nivel de servicio.
audIT responde solidariamente por todos ellos \art*{73.3}{38}.

% ============================================================================
\section{Gestión del conocimiento y continuidad}\label{sec:12-gestion-del-conocimiento-y-conti}
% ============================================================================

El \art{76.4}{39} pide que la salida de cualquier persona no comprometa el proyecto. La regla de audIT
es que todo conocimiento crítico lo tengan al menos dos personas, con la revisión cruzada que describe el
Subdocumento 1. Cada rol del equipo clave tiene un suplente designado dentro de la planta. El Jefe de
Proyecto tiene como suplente a uno de los tres analistas de la oficina de proyecto. El Arquitecto de
Solución y el Líder de Desarrollo se cubren entre sí. El Encargado de Seguridad tiene como suplente al
Líder de Operación, que ejecuta con él la identidad y el cifrado. El Líder de Datos tiene al ingeniero
de datos, la Líder de Calidad al ingeniero de pruebas y automatización y el Líder de Operación a uno de
los ingenieros de la operación de la plataforma. Las decisiones quedan en el registro de arquitectura,
los procedimientos en los libros de operación y las soluciones en la base de conocimiento (Subdocumento
10), de modo que lo que se sabe no depende de una sola memoria.

Las 16 personas de refuerzo entran en pareja con un profesional de planta que ya conoce su frente, y su
primer trabajo se revisa antes de integrarse. Cada una recibe su cuenta del repositorio y su identidad
en Entra ID con acceso condicional y doble factor, y el día en que deja el proyecto se le retiran los
accesos. Así el peak de 37 personas no deja conocimiento en manos de quien se va en el mes 18.

Un reemplazo del equipo clave necesita la autorización previa y escrita del mandante y recae en un
profesional de perfil igual o superior, con quince días hábiles de traslape y sin costo para Curimón
\art*{76.2}{39}. La rotación no puede pasar del 30 \% del equipo clave en doce meses \art*{76.3}{39}, que
con seis personas significa no más de una salida al año. Para retener al equipo, audIT aplica durante
todo el contrato el programa de entrenamiento y certificación del Subdocumento 1.

El conocimiento también pasa al mandante. Las personas administradoras y el equipo técnico de Curimón se
evalúan y certifican al cierre de cada marcha blanca \transv*{RT-22.05}{38}. Los libros de operación los
prueba el personal de Curimón antes de recibirlos \transv*{RT-21.04}{35}, y después del paso a
producción de la Etapa 2 audIT acompaña al equipo técnico del mandante durante al menos seis meses
\transv*{RT-22.07}{38}. Durante la operación siguen las dos jornadas anuales de actualización del
Subdocumento 10.
